\section{Validation and uncertainty}
\label{app:validation}

\subsection{What has been validated}

The quantitative checks that this study actually performs are:

\begin{enumerate}
    \item \textbf{Internal consistency of the random limit}: the two random-growth variants (independent random percolation and cumulative density increment) give gel-points agreeing to $\Delta p = 0.0002$ (main text).
    \item \textbf{Geometric anchor}: at the random limit the dynamical gel-point lies within one percent of the complementary void-percolation threshold $1 - \pc = 0.6884$.
    \item \textbf{Adjacency-rule robustness}: the 6-neighbour and 26-neighbour many-nuclei variants agree to $|\Delta\pcprime| = 0.0005$.
    \item \textbf{Replication}: the $\kappa$-family gel-points are means over $N_s = 3$ independent seeds, and the reported uncertainties are the corresponding standard deviations.
\end{enumerate}

\subsection{What has not been done}

For clarity, the following were \emph{not} carried out, and no results are claimed for them: cross-validation, bootstrap resampling, hypothesis tests ($t$-tests, ANOVA), effect-size calculations, and formal uncertainty propagation. In particular, no prediction-accuracy percentage is reported. The earlier draft of this appendix quoted such figures; they were not computed from the data and have been removed.

\subsection{Sources of uncertainty}

The dominant uncertainty on a quoted gel-point is the resolution of the occupation grid together with the replicate scatter. Where a sigmoid fit is used (Appendix~\ref{app:mathematical}), the formal fit uncertainty is reported alongside the value; for the many-nuclei morphology that formal uncertainty is large unless the curve is padded, and this is stated explicitly rather than suppressed. The GSER spectra are single-replicate per condition, so their moduli are less tightly constrained than the gel-points.
